
% Copyright (c)  2007-2012  Tim Niemueller, RWTH Aachen University
%
% Created: Mon Aug 13 16:55:22 2012

%aspectratio=169
\documentclass[compress,xcolor=table,nominiframes,notitlehead,noframetitlebg,addurl]{beamer}
% add handout to optional args for handout version

%\usepackage{hyperref}
\usepackage{beamerthemesplit}
\usepackage[utf8]{inputenc}
\usepackage[T1]{fontenc}
\usepackage[german]{babel}
\usepackage{german}
\usepackage{graphicx}
\usepackage{subfigure}
\usepackage{epsfig}
\usepackage{tabularx}
\usepackage{latexsym}
\usepackage{url}
\usepackage{tikz}
\usepackage{multimedia}
\usepackage{xcolor}
\usepackage{listings}
\usepackage{verbatim}
\usepackage{colortbl}
\usepackage{psfrag}
\usepackage{xifthen}
\usepackage[absolute,overlay]{textpos}
\usepackage{palatino}
\usepackage{enumitem}

\graphicspath{{../},{./images/},{./images/},{./figures/}}

%% Mulberry Color for highlighting
\definecolor{Mulberry}{cmyk}{0.34,0.90,0,0.02}
\definecolor{Lavender}{cmyk}{0,0.48,0,0}
\definecolor{Melon}{cmyk}{0,0.46,0.50,0}
\definecolor{Peach}{cmyk}{0,0.50,0.70,0}
\definecolor{RedOrange}{cmyk}{0,0.77,0.87,0}
\definecolor{BrickRed}{cmyk}{0,0.89,0.94,0.28}
\definecolor{Mahogany}{cmyk}{0,0.85,0.87,0.35}
\definecolor{BurntOrange}{cmyk}{0,0.51,1,0}
\definecolor{BitterSweet}{cmyk}{0,0.75,1,0.24}
\definecolor{RosBlue}{rgb}{0.19,0.25,0.38}


\mode<presentation>
{
  \usetheme{FawkesMinimal}

  % to hide nav bar uncomment this line
  \setbeamertemplate{navigation symbols}{}

  %\setbeamercovered{transparent}
  \setbeamercovered{%
    invisible,
    again covered={\opaqueness<1->{40}}
  }

  \setbeamerfont{author}{size=\footnotesize,family=\rmfamily,parent=structure}
}

% Usage notes for handout version:
% Compile the beamer version immediately before you build the handout version,
% otherwise page numbers etc. will be wrong! The .aux files are *not* updated
% in handout mode, see PGF Manual for details why this is necessary.
% Comment out below one of the two pgfuselayout lines for either 2 or 4 slides
% per page. To have a very light grey background uncomment the background canvas
% color line. The logical page options are used to draw borders around each
% slide.
\mode<handout>
{
  \usetheme{Fawkes}

  % to hide nav bar uncomment this line
  \setbeamertemplate{navigation symbols}{}

  %\setbeamercovered{transparent}
  \setbeamercovered{%
    again covered={\opaqueness<1->{40}}
  }

  % Very slight grey background, can be used instead of borders
  %\setbeamercolor{background canvas}{bg=black!5}

  \usepackage{pgfpages}
  \pgfpagesuselayout{4 on 1}[a4paper,border shrink=5mm,landscape]
  %\pgfpagesuselayout{2 on 1}[a4paper,border shrink=5mm]

  \pgfpageslogicalpageoptions{1}{border code=\pgfstroke}
  \pgfpageslogicalpageoptions{2}{border code=\pgfstroke}
  \pgfpageslogicalpageoptions{3}{border code=\pgfstroke}
  \pgfpageslogicalpageoptions{4}{border code=\pgfstroke}
  \nofiles
}


% Declare layers
\pgfdeclarelayer{background}
\pgfsetlayers{background,main} 

% Load PGF libraries
\usetikzlibrary{patterns}
\usetikzlibrary{arrows}
\usetikzlibrary{topaths}
\usetikzlibrary{snakes}
\usetikzlibrary{calc}
\usetikzlibrary{positioning}
\usetikzlibrary{shadows}
\usetikzlibrary{shapes.multipart}
%% \usetikzlibrary{arrows.meta}
\usetikzlibrary{fit,arrows,decorations.pathreplacing,decorations.text}


% set lengths for textpos package
\setlength{\TPHorizModule}{10mm}
\setlength{\TPVertModule}{\TPHorizModule}
\textblockorigin{8mm}{16mm} % start everything near the top-left corner
\setbeamercolor{textblock color}{fg=blue!50,bg=white}

\institute{%
  %\vspace{1cm}
  \begin{minipage}{\textwidth}\centering
  \includegraphics[height=0.6cm]{common/images/logos/rwth-logo}
  \quad\quad
  \includegraphics[height=0.6cm]{common/images/logos/kbsg}
  \end{minipage}
}

\setlist[itemize]{leftmargin=12pt,label=\usebeamerfont*{itemize item}%
  \usebeamercolor[fg]{itemize item}
  \usebeamertemplate{itemize item}}
\setlist[description]{leftmargin=4pt,itemindent=0pt,font=\color{FawkesOrange}}
\setenumerate[1]{%
  leftmargin=12pt,%
  label=\protect\usebeamerfont{enumerate item}%
        \protect\usebeamercolor[fg]{enumerate item}%
        \insertenumlabel.}

%numbers=left, numberstyle=\tiny, stepnumber=2, numbersep=5pt
\lstset{language=[GNU]C++,
        basicstyle=\small,
        escapeinside={/*(*/}{/*)*/},
        breaklines=true,
        showstringspaces=false
        }

% Uncomment the listings styles you need
\lstdefinelanguage{Lua}
{
  morekeywords={and,break,do,else,elseif,end,false,for,function,
                if,in,local,nil,not,or,repeat,return,then,true,until,while},
  sensitive=true,
  morecomment=[l]{--},
  morecomment=[s]{--[[}{--]]},
  morestring=[b]{"},
  morestring=[s]{[==[}{]==]},
}

% default style
\lstdefinestyle{Lua}
{
  language=Lua,
  basicstyle=\ttfamily,
  breaklines=true,
  showstringspaces=false,
  %keywordstyle=\bfseries,
  keywordstyle=\color{Mulberry},
  %frame=lines,
  %belowcaptionskip=8pt,
  emphstyle=\itshape,
  %numbers=left,
  stepnumber=1,
  %backgroundcolor=\color{blue!10},
  rulecolor=\color{blue!50},
  fillcolor=\color{blue!20},
  %framexleftmargin=18pt,
  %xleftmargin=18pt,
  stringstyle=\color{BitterSweet},
  %stringstyle=\color{BrickRed},
  commentstyle=\color{BrickRed},
  escapechar=\%
  % emph={getup, servo, depends_skills},
  %emphstyle=\underbar,
  %numbers=left,
  %stepnumber=1,
  %%stringstyle=\ttfamily, % typewriter type for strings
}
\lstdefinestyle{SmallLua}{
  style=Lua,
  basicstyle=\ttfamily\footnotesize,
  numbersep=6pt,
}
\lstdefinestyle{ReallySmallLua}{
  style=Lua,
  basicstyle=\ttfamily\tiny,
  numbersep=5pt,
}

\lstdefinelanguage{Yaml}
{
  keywords={true,false,null,y,n},
  sensitive=false,
  comment=[l]{\#},
  morecomment=[s]{/*}{*/},
  morestring=[b]',
  morestring=[b]",
}

% default style
\lstdefinestyle{Yaml}
{
  language=Yaml,
  basicstyle=\YAMLkeystyle,
  breaklines=true,
  showstringspaces=false,
  %keywordstyle=\bfseries,
  keywordstyle=\color{Mulberry},
  %frame=lines,
  %belowcaptionskip=8pt,
  emphstyle=\itshape,
  %numbers=left,
  stepnumber=1,
  %backgroundcolor=\color{blue!10},
  rulecolor=\color{blue!50},
  fillcolor=\color{blue!20},
  %framexleftmargin=18pt,
  %xleftmargin=18pt,
  stringstyle=\color{BitterSweet},
  %stringstyle=\color{BrickRed},
  commentstyle=\color{BrickRed},
  escapechar=\%
  % emph={getup, servo, depends_skills},
  %emphstyle=\underbar,
  %numbers=left,
  %stepnumber=1,
  %%stringstyle=\ttfamily, % typewriter type for strings
}
\lstdefinestyle{SmallYaml}{
  style=Yaml,
  basicstyle=\ttfamily\footnotesize,
  numbersep=6pt,
}
\lstdefinestyle{ReallySmallYaml}{
  style=Yaml,
  basicstyle=\ttfamily\tiny,
  numbersep=5pt,
}

%\lstdefinelanguage{OpenPRS}{
  keywordsprefix=\$,
  %keywordsprefix=\$,
  alsoletter={?!=-<>*\$^~},
  keywordstyle=\color{Mulberry!80!black}\bfseries,
  keywords=[2]{:invocation, :context, :call, :body, :action,
    :setting, :documentation, :effects, defop},
  keywordstyle=[2]\color{BrickRed!70!blue}\bfseries,
  keywords=[3]{?, !, ^, ~, =>, ~>, =, if, else, elseif, do, while, //},
  keywordstyle=[3]\color{darkgray}\bfseries,
  keywords=[4]{start-critical-section, end-critical-section, equal,
    print, printf},%
  keywordstyle=[4]\color{gray}\bfseries,
  %identifierstyle=\color{black},
  sensitive=false,
  comment=[l]{;},
  commentstyle=\color{purple}\ttfamily,
  stringstyle=\color{red}\ttfamily,
  morestring=[b]"
}


\lstdefinestyle{OpenPRS}
{
  language=OpenPRS,
  basicstyle=\footnotesize\ttfamily\vspace{0.2cm},
  breaklines=true,
  showstringspaces=false,
  %keywordstyle=\bfseries,
  %keywordstyle=\color{Mulberry},
  frame=lines,
  belowcaptionskip=8pt,
  emphstyle=\itshape,
  numbers=left,
  stepnumber=1,
  backgroundcolor=\color{blue!10},
  rulecolor=\color{blue!50},
  fillcolor=\color{blue!20},
  framexleftmargin=16pt,
  xleftmargin=16pt,
  %stringstyle=\color{BitterSweet},
  stringstyle=\color{BrickRed},
  commentstyle=\color{BrickRed},
  escapechar=\%
  % emph={getup, servo, depends_skills},
  %emphstyle=\underbar,
  %numbers=left,
  %stepnumber=1,
  %%stringstyle=\ttfamily, % typewriter type for strings
}

\lstdefinestyle{SmallOpenPRS}{
  style=OpenPRS,
  basicstyle=\ttfamily\footnotesize,
  numbersep=6pt,
}
\lstdefinestyle{ReallySmallOpenPRS}{
  style=OpenPRS,
  basicstyle=\ttfamily\scriptsize,
  numbersep=5pt,
}

\lstdefinestyle{ReallySmallOpenPRSNoFrame}{
  style=OpenPRS,
  basicstyle=\ttfamily\scriptsize,
  numbersep=5pt,
  frame=none,
  backgroundcolor=\color{white},
  framextopmargin=0pt,
  framexbottommargin=0pt
}

\lstdefinestyle{SuperSmallOpenPRSNoFrame}{
  style=OpenPRS,
  basicstyle=\ttfamily\fontsize{10pt}{10pt}\selectfont,
  numbers=none,
  frame=none,
  backgroundcolor=\color{white},
  framextopmargin=0pt,
  framexbottommargin=0pt,
  framexleftmargin=-2pt, xleftmargin=-2pt,
}

\lstdefinestyle{SmallOpenPRSNoFrame}{
  style=OpenPRS,
  basicstyle=\ttfamily\footnotesize,
  numbersep=5pt,
  frame=none,
  backgroundcolor=\color{white},
  framextopmargin=0pt,
  framexbottommargin=0pt
}

\lstdefinelanguage{CLIPS}{
  keywordsprefix=?,
  %keywordsprefix=\$,
  alsoletter={?=-<>*\$},
  keywordstyle=\color{Mulberry!80!black}\bfseries,
  keywords=[2]{deffunction, deftemplate, defrule, deffacts, defgeneric,
    defmodule, defadvice, defglobal, defmethod, definstance, defclass},
  keywordstyle=[2]\color{BrickRed!70!blue}\bfseries,
  keywords=[3]{slot, multislot, type, default, default-dynamic,
                      extends, crlf, range, nil, if, then, else, while,
                      not, or, switch, case, and, reset,
                      assert, test, declare, salience, return, bind, modify,
                      retract, explicit, unique, node-index-hash,
                      halt, printout, =>, <-},
  keywordstyle=[3]\color{darkgray}\bfseries,
  keywords=[4]{subsetp,progn, progn$, not, node-index-hash, create$,
    append$, length$, printout},%
  keywordstyle=[4]\color{gray}\bfseries,
  %identifierstyle=\color{black},
  sensitive=false,
  comment=[l]{;},
  commentstyle=\color{purple}\ttfamily,
  stringstyle=\color{red}\ttfamily,
  morestring=[b]"
}


\lstdefinestyle{CLIPS}
{
  language=CLIPS,
  basicstyle=\footnotesize\ttfamily\vspace{0.2cm},
  breaklines=true,
  showstringspaces=false,
  %keywordstyle=\bfseries,
  %keywordstyle=\color{Mulberry},
  frame=lines,
  belowcaptionskip=-3pt,
  emphstyle=\itshape,
  numbers=left,
  stepnumber=1,
  backgroundcolor=\color{gray!10},
  rulecolor=\color{gray!80},
  fillcolor=\color{gray!10},
  framexleftmargin=16pt,
  xleftmargin=16pt,
  %stringstyle=\color{BitterSweet},
  stringstyle=\color{BrickRed},
  commentstyle=\color{BrickRed},
  escapechar=\%,
  % emph={getup, servo, depends_skills},
  %emphstyle=\underbar,
  %numbers=left,
  %stepnumber=1,
  %%stringstyle=\ttfamily, % typewriter type for strings
  %float,
  captionpos=b
}

\lstdefinestyle{SmallCLIPS}{
  style=CLIPS,
  basicstyle=\ttfamily\footnotesize,
  numbersep=6pt,
}
\lstdefinestyle{ReallySmallCLIPS}{
  style=CLIPS,
  basicstyle=\ttfamily\scriptsize,
  numbersep=5pt,
}

\lstdefinestyle{ReallySmallCLIPSNoFrame}{
  style=CLIPS,
  basicstyle=\ttfamily\scriptsize,
  numbersep=5pt,
  frame=none,
  backgroundcolor=\color{white},
  framextopmargin=0pt,
  framexbottommargin=0pt
}

\lstdefinestyle{SuperSmallCLIPSNoFrame}{
  style=CLIPS,
  basicstyle=\ttfamily\fontsize{10pt}{10pt}\selectfont,
  numbers=none,
  frame=none,
  backgroundcolor=\color{white},
  framextopmargin=0pt,
  framexbottommargin=0pt,
  framexleftmargin=-2pt, xleftmargin=-2pt,
}

\lstdefinestyle{SmallCLIPSNoFrame}{
  style=CLIPS,
  basicstyle=\ttfamily\footnotesize,
  numbersep=5pt,
  frame=none,
  backgroundcolor=\color{white},
  framextopmargin=0pt,
  framexbottommargin=0pt
}

%\lstdefinelanguage{JavaScript}{
  keywords={typeof, new, true, false, catch, function, return, null, catch, switch, var, if, in, while, do, else, case, break},
  keywordstyle=\color{blue}\bfseries,
  ndkeywords={class, export, boolean, throw, implements, import}, %, this
  ndkeywordstyle=\color{darkgray}\bfseries,
  identifierstyle=\color{black},
  sensitive=false,
  comment=[l]{//},
  morecomment=[s]{/*}{*/},
  commentstyle=\color{purple}\ttfamily,
  stringstyle=\color{red}\ttfamily,
  morestring=[b]',
  morestring=[b]"
}

\lstdefinestyle{JSON}
{
  language=JavaScript,
  morekeywords={interface,field,message,comment},
  basicstyle=\footnotesize\ttfamily\vspace{0.2cm},
  breaklines=true,
  showstringspaces=false,
  %keywordstyle=\bfseries,
  keywordstyle=\color{Mulberry},
  frame=lines,
  belowcaptionskip=8pt,
  emphstyle=\itshape,
  numbers=left,
  stepnumber=1,
  backgroundcolor=\color{blue!10},
  rulecolor=\color{blue!50},
  fillcolor=\color{blue!20},
  framexleftmargin=16pt,
  xleftmargin=16pt,
  %stringstyle=\color{BitterSweet},
  stringstyle=\color{BrickRed},
  commentstyle=\color{BrickRed},
  escapechar=\%,
  % emph={getup, servo, depends_skills},
  %emphstyle=\underbar,
  %numbers=left,
  %stepnumber=1,
  %%stringstyle=\ttfamily, % typewriter type for strings
  captionpos=b
}

\lstdefinestyle{SmallJSON}{
  style=JSON,
  basicstyle=\ttfamily\footnotesize,
  numbersep=6pt,
}
\lstdefinestyle{ReallySmallJSON}{
  style=JSON,
  basicstyle=\ttfamily\tiny,
  numbersep=5pt,
}

\definecolor{dkgreen}{rgb}{0,0.6,0}
\definecolor{gray}{rgb}{0.5,0.5,0.5}
\definecolor{mauve}{rgb}{0.58,0,0.82}
\definecolor{gray}{rgb}{0.4,0.4,0.4}
\definecolor{darkblue}{rgb}{0.0,0.0,0.6}
\definecolor{lightblue}{rgb}{0.0,0.0,0.9}
\definecolor{cyan}{rgb}{0.0,0.6,0.6}
\definecolor{darkred}{rgb}{0.6,0.0,0.0}


\lstdefinelanguage{XML}
{
  basicstyle=\ttfamily\footnotesize,
  columns=fullflexible,
  showstringspaces=false,
  numbers=left,                   % where to put the line-numbers
  numberstyle=\tiny\color{gray},  % the style that is used for the line-numbers
  stepnumber=1,
  numbersep=5pt,                  % how far the line-numbers are from the code
  backgroundcolor=\color{white},      % choose the background color. You must add \usepackage{color}
  showspaces=false,               % show spaces adding particular underscores
  showstringspaces=false,         % underline spaces within strings
  showtabs=false,                 % show tabs within strings adding particular underscores
  frame=none,                   % adds a frame around the code
  rulecolor=\color{black},        % if not set, the frame-color may be changed on line-breaks within not-black text (e.g. commens (green here))
  tabsize=2,                      % sets default tabsize to 2 spaces
  captionpos=b,                   % sets the caption-position to bottom
  breaklines=true,                % sets automatic line breaking
  breakatwhitespace=false,        % sets if automatic breaks should only happen at whitespace
  title=\lstname,                   % show the filename of files included with \lstinputlisting;
                                  % also try caption instead of title  
  commentstyle=\color{gray}\upshape,
  morestring=[s][\color{mauve}]{"}{"},
  morestring=[s][\color{black}]{>}{<},
  morecomment=[s]{<?}{?>},
  morecomment=[s][\color{dkgreen}]{<!--}{-->},
  stringstyle=\color{black},
  identifierstyle=\color{lightblue},
  keywordstyle=\color{red},
  morekeywords={xmlns,xsi,noNamespaceSchemaLocation,type,id,x,y,source,target,version,tool,transRef,roleRef,objective,eventually},
  basicstyle=\ttfamily\scriptsize,
  backgroundcolor=\color{blue!10},
  numbers=none,
  numbersep=5pt
}


% Default style
\lstset{style=Lua}

% Hyphenation of words with hyphen
\def\hyph{-\penalty0\hskip0pt\relax}

% define an anchor in the frame
\newcommand{\tikzref}[1]{%
  \tikz[remember picture]{%
    \coordinate (#1) at (0,0.5ex);%
  }%
}%


%%% Local Variables: 
%%% mode: latex
%%% TeX-master: "aaai-spring2013-clips-agent"
%%% End: 
